% ===========================================================
%  第 2 章 余子式和代数余子式的计算 —— 高等代数【深化】拓展集（精选）
%  来源：主控自撰 + 北大《高等代数》第五版 §4 矩阵的逆
%        + 樊启斌《高等代数典型问题与方法》§2.3 代数余子式问题
%  筛选依据：AGENTS.md §2 决定二「只补能加深理解、反哺解题的深化知识」
% ===========================================================
%
%  【保留 12 块】
%   深化一 伴随矩阵与可逆性：伴随矩阵一个恒等式定全局(自撰) / r(A*) 三分(自撰) /
%             逆矩阵的定义与唯一性(北大) / 矩阵可逆的充要条件与逆矩阵公式(北大) /
%             只验一边即可判定互逆(北大) / 逆矩阵的运算性质(北大)
%   深化二 代数余子式的和：代数余子式的和怎么求(自撰) / |A*|=|A|^{n-1}及推论(自撰) /
%             余子式之和与 |A+xJ|(樊) / 行和列和全为0⇒代数余子式全相等(樊) /
%             A_ij=a_ij ⇒ A*=A^T 型(樊)
%   深化三 与方程组衔接：用逆矩阵推导克拉默法则(北大)
%
%  【删除及理由】
%   · 北大「伴随矩阵基本恒等式 AA*=|A|E」→ 自撰块「伴随矩阵：一个恒等式定全局」已给
%     完整证明（含非对角元"异行代数余子式之和为 0"的来源），取一即可。
%   · 北大「伴随矩阵的定义：代数余子式的转置排布(定义9)」→ 已并入自撰块 1，
%     且自撰版额外强调"转置"这一高频失分点。
%   · 樊「余子式的两个基本计算式（升阶法与降阶法）」→ 与自撰块「代数余子式的和怎么求」重复。
%   · 樊「余子式之和、代数余子式之和的通用求法（加全1行/列）」→ 同上，重复。
%   · 樊「代数余子式的线性组合：行替换公式与列替换公式」→ 同上，重复。
%   · 樊「对全排列求和的抽象行列式必为 0」→ 竞赛级结论，考纲无接口。
%
%  【来源疑点已由主控核实并排除】
%   转录代理举报樊书印 107–108「|A+J|=|A|+|A_3|」有误，并给出反例 A=[[2,1],[1,1]]。
%   主控数值复核：|A+J|=2（代理误算为 6），而 |A|+ΣΣA_ij=1+1=2，**恰好印证**书中公式；
%   且核心恒等式 |A+xJ|=|A|+x·ΣΣA_ij 在 n=2,3,4、x∈{0,±1,1.5,3} 上全部成立。
%   → **不采信该举报**。所收结论均已独立数值验证。该例题本就只收结论、不抄演算。
% ===========================================================

\fsec{深化一　伴随矩阵与矩阵可逆性}

\begin{fdeep}{伴随矩阵：一个恒等式定全局}
\textbf{定义}：设 $A=(a_{ij})_{n\times n}$，$A_{ij}$ 为 $a_{ij}$ 的代数余子式，则
\fml{A^{*}=\begin{pmatrix}A_{11}&A_{21}&\cdots&A_{n1}\\ A_{12}&A_{22}&\cdots&A_{n2}\\ \vdots&\vdots&&\vdots\\ A_{1n}&A_{2n}&\cdots&A_{nn}\end{pmatrix}}
\textbf{注意转置}：$A^{*}$ 的第 $i$ \textbf{行}是 $A$ 第 $i$ \textbf{列}元素的代数余子式。这一条错位是最常见的失分点。

\medskip
\textbf{唯一需要记住的恒等式}：
\fml{AA^{*}=A^{*}A=|A|E}
\textbf{它从哪来}：$(AA^{*})_{ij}=\sum_{k=1}^{n}a_{ik}A_{jk}$，即"第 $i$ 行元素乘第 $j$ 行元素的代数余子式"。
由行列式展开定理，
\fml{\sum_{k=1}^{n}a_{ik}A_{jk}=\begin{cases}|A|,& i=j\\ 0,& i\ne j\end{cases}}
——$i=j$ 时是"按第 $i$ 行展开"，$i\ne j$ 时相当于把第 $i$ 行换成第 $j$ 行（两行相同，行列式为 $0$）。
所以 $AA^{*}=|A|E$ 不是新结论，而是\textbf{展开定理的直接改写}。

\medskip
\textbf{由它一网打尽全部常用结论}（$n\ge2$ 时）：
\begin{itemize}[leftmargin=2.2em,itemsep=0.22em]
\item $|A|\ne0$ 时，$A^{-1}=\dfrac{1}{|A|}A^{*}$ —— 这就是\textbf{用伴随矩阵求逆}的依据；
\item $|A^{*}|=|A|^{\,n-1}$ —— 对 $AA^{*}=|A|E$ 两边取行列式，$|A|\,|A^{*}|=|A|^{n}$，
      再对 $|A|=0$ 的情形作连续性/按定义单独论证即得；
\item $(A^{*})^{*}=|A|^{\,n-2}A$（$n\ge3$）；
\item $(kA)^{*}=k^{\,n-1}A^{*}$，$(A^{\mathrm{T}})^{*}=(A^{*})^{\mathrm{T}}$；
\item $A^{*}=|A|A^{-1}$ 说明\textbf{$A^{*}$ 与 $A$ 是"同类"的}：$A$ 对称则 $A^{*}$ 对称，$A$ 可逆则 $A^{*}$ 可逆。
\end{itemize}
【反哺点】考纲"矩阵"一节要求"理解伴随矩阵的概念，会用伴随矩阵求逆矩阵"。
把这一串结论全部挂在 $AA^{*}=|A|E$ 这一个等式上，就不必逐条背——
考场上先写出 $AA^{*}=|A|E$，再问"我要的是哪一项"，绝大多数伴随矩阵题都能现场推出来。
\end{fdeep}

\begin{fdeep}{$r(A^{*})$ 的三分判别：$n$ / $1$ / $0$}
设 $A$ 为 $n$ 阶矩阵（$n\ge2$），则
\fml{r(A^{*})=\begin{cases}n,& r(A)=n\\ 1,& r(A)=n-1\\ 0,& r(A)\le n-2\end{cases}}
\textbf{三段证明思路}：
\begin{enumerate}[leftmargin=2.2em,itemsep=0.25em]
\item $r(A)=n$（即 $|A|\ne0$）：由 $|A^{*}|=|A|^{\,n-1}\ne0$，故 $A^{*}$ 可逆，$r(A^{*})=n$；
\item $r(A)=n-1$：此时 $A$ 存在一个 $n-1$ 阶非零子式，故\textbf{至少有一个 $A_{ij}\ne0$}，即 $A^{*}\ne O$，$r(A^{*})\ge1$；
      又由 $AA^{*}=|A|E=O$ 知 $A^{*}$ 的每一列都是 $Ax=0$ 的解向量。
      而 $r(A)=n-1$ 时 $Ax=0$ 的基础解系只含 $n-r(A)=1$ 个向量，
      故 $A^{*}$ 的各列都与同一个向量成比例，$r(A^{*})\le1$。合起来 $r(A^{*})=1$；
\item $r(A)\le n-2$：此时 $A$ 的所有 $n-1$ 阶子式全为 $0$，故全部 $A_{ij}=0$，即 $A^{*}=O$，$r(A^{*})=0$。
\end{enumerate}
\textbf{第 ② 段的思路最值得记}：它把"伴随矩阵的秩"转化为\textbf{"解空间的维数只有 1"}\allowbreak——
这正好复用第 5 章"基础解系含 $n-r$ 个向量"的结论。

【反哺点】这是考研选择题的高频送分点，标准问法是"已知 $r(A)=n-1$，求 $r(A^{*})$"。
另外它给出一个反向用法：\textbf{由 $r(A^{*})$ 反推 $r(A)$}——见到 $A^{*}=O$ 就意味着 $r(A)\le n-2$。
\end{fdeep}

\begin{fdeep}{逆矩阵的定义与唯一性}
首先我们指出，由于矩阵的乘法规则，只有方阵才能满足 $AB=BA=E$（读者自己证明）.其次，对于任意的矩阵 $A$，适合等式 $AB=BA=E$ 的矩阵 $B$ 是唯一的（如果有的话）.事实上，假设 $B_1,B_2$ 是两个适合它的矩阵，就有
\fml{B_1=B_1E=B_1(AB_2)=(B_1A)B_2=EB_2=B_2.}

\textbf{定义 8}\quad 如果矩阵 $B$ 适合 $AB=BA=E$，那么 $B$ 就称为 $A$ 的逆矩阵，记作 $A^{-1}$.下面要解决的问题是：在什么条件下矩阵 $A$ 是可逆的？如果 $A$ 可逆，怎样求 $A^{-1}$？

【反哺点】服务考纲「矩阵」中"理解逆矩阵的概念、掌握逆矩阵的性质"：唯一性的证明只用到"$E$ 是乘法单位元 $+$ 结合律"，是"证唯一性"这类小题的标准四步；同时点明逆矩阵只对方阵才有意义，可避免对非方阵误求逆。
\end{fdeep}

\begin{fdeep}{矩阵可逆的充要条件与逆矩阵公式（定理 3）}
如果 $d=|A|\ne 0$，那么由 (2) 得
\fml{A\left(\frac{1}{d}A^{*}\right)=\left(\frac{1}{d}A^{*}\right)A=E.\qquad (3)}

\textbf{定理 3}\quad 矩阵 $A$ 是可逆的充分必要条件是 $A$ 非退化，而
\fml{A^{-1}=\frac{1}{d}A^{*}\qquad (d=|A|\ne 0).}

证明思路（充分性）：当 $d=|A|\ne 0$ 时，由 (3) 可知 $A$ 可逆，且
\fml{A^{-1}=\frac{1}{d}A^{*}.\qquad (4)}

证明思路（必要性）：反过来，如果 $A$ 可逆，那么有 $A^{-1}$ 使 $AA^{-1}=E$.两边取行列式，得
\fml{|A|\,|A^{-1}|=|E|=1,\qquad (5)}
因而 $|A|\ne 0$，即 $A$ 非退化.

【反哺点】服务考纲「矩阵」"会用伴随矩阵求逆矩阵"：$A^{-1}=\dfrac{1}{|A|}A^{*}$ 就是考纲原话的公式形态，二阶情形 $\begin{pmatrix}a&b\\c&d\end{pmatrix}^{-1}=\dfrac{1}{ad-bc}\begin{pmatrix}d&-b\\-c&a\end{pmatrix}$ 由它直接得到.证明演示了两条通用手法——"先凑出 $AX=E$ 即可断言可逆"与"等式两边取行列式".而"$|A|\ne0\Leftrightarrow A$ 可逆"是贯穿方程组、秩、特征值讨论的总开关.
\end{fdeep}

\begin{fdeep}{只验一边即可判定互逆，以及 $|A^{-1}|=|A|^{-1}$}
根据定理 3 容易看出，对于 $n$ 阶方阵 $A,B$，如果
\fml{AB=E,}
那么 $A,B$ 就都是可逆的并且它们互为逆矩阵.

由 (5) 可以看出，如果 $d=|A|\ne 0$，那么
\fml{|A^{-1}|=d^{-1}.}

（另注：定理 3 不但给出了矩阵可逆的条件，同时也给出了求逆矩阵的公式 (4).按这个公式来求逆矩阵，计算量一般是很大的.）

【反哺点】服务考纲「矩阵」：判定 $A,B$ 互逆只需验证一边 $AB=E$，不必再验 $BA=E$——选项题与证明题里省掉一半工作量；$|A^{-1}|=|A|^{-1}$ 常与 $|A^{*}|=|A|^{n-1}$、$|kA|=k^{n}|A|$ 串联出现，是行列式数值计算的固定链条；末句提示伴随公式适合理论推导而非大阶数手算（大阶数应改用初等行变换法），可避免考场上硬算伴随矩阵。
\end{fdeep}

\begin{fdeep}{逆矩阵的运算性质：转置与乘积（推论）}
\textbf{推论}\quad 如果矩阵 $A,B$ 可逆，那么 $A^{\mathrm{T}}$ 与 $AB$ 也可逆，且
\fml{(A^{\mathrm{T}})^{-1}=(A^{-1})^{\mathrm{T}},\qquad (AB)^{-1}=B^{-1}A^{-1}.}

证明思路：由定理即得推论的前一半，现在来证后一半.由
\fml{AA^{-1}=A^{-1}A=E}
两边取转置，有
\fml{(A^{-1})^{\mathrm{T}}A^{\mathrm{T}}=A^{\mathrm{T}}(A^{-1})^{\mathrm{T}}=E^{\mathrm{T}}=E,}
因此 $(A^{\mathrm{T}})^{-1}=(A^{-1})^{\mathrm{T}}$.由
\fml{(AB)(B^{-1}A^{-1})=(B^{-1}A^{-1})(AB)=E}
即得 $(AB)^{-1}=B^{-1}A^{-1}$.

【反哺点】服务考纲「矩阵」"掌握逆矩阵的性质"：证明给出两条通用手段——"等式两边取转置"与"直接验证乘积等于 $E$"（配合上一条的"只验一边"）.$(AB)^{-1}=B^{-1}A^{-1}$ 的顺序颠倒与 $(AB)^{\mathrm{T}}=B^{\mathrm{T}}A^{\mathrm{T}}$ 完全同构，是求分块矩阵逆、解矩阵方程 $AXB=C$（得 $X=A^{-1}CB^{-1}$）时的关键，绝不能写成 $A^{-1}B^{-1}$.
\end{fdeep}

\fsec{深化二　代数余子式的和与行列式构造}

\begin{fdeep}{代数余子式的和怎么求：把"配系数"翻译成"构造行列式"}
本讲的核心手法是：\textbf{代数余子式前的系数，就是替换掉的那一行（列）的元素}。
\fml{\sum_{k=1}^{n}k_{k}A_{ik}=\begin{vmatrix}
a_{11}&\cdots&a_{1n}\\
\vdots&&\vdots\\
k_{1}&\cdots&k_{n}&\leftarrow\text{第 }i\text{ 行被换成 }k_{1},\dots,k_{n}\\
\vdots&&\vdots\\
a_{n1}&\cdots&a_{nn}
\end{vmatrix}}
也就是说：\textbf{$A_{i1},A_{i2},\dots,A_{in}$ 只由 $A$ 中\emph{去掉第 $i$ 行}后的部分决定}，
与第 $i$ 行元素本身无关。既然无关，就可以\textbf{自由地把第 $i$ 行换成任何想要的行}——
只要那个新行列式好算，就算出了这组代数余子式的线性组合。

\medskip
\textbf{由此得到一套"求余子式之和"的固定套路}：
\begin{enumerate}[leftmargin=2.2em,itemsep=0.25em]
\item 要把 $\sum A_{i1}$（第 $i$ 行代数余子式之和）求出来，
      就在 $\sum_k a_{ik}A_{ik}=|A|$ 里\textbf{把 $a_{ik}$ 全部换成 $1$}，得
      \fml{\sum_{k=1}^{n}A_{ik}=\begin{vmatrix}\cdots\\ 1&1&\cdots&1\\ \cdots\end{vmatrix}}
      即"把第 $i$ 行整行换成 $1$"之后的行列式；
\item 要求 $\sum_{i,j}A_{ij}$（\textbf{全部}代数余子式之和），就把它看成
      $\sum_i\left(\sum_j A_{ij}\right)$：对每个 $i$ 都把第 $i$ 行换成全 $1$，再全部相加。
      更快的一步到位写法是左乘全 $1$ 矩阵：
      \fml{\sum_{i=1}^{n}\sum_{j=1}^{n}A_{ij}=\mathbf{1}^{\mathrm{T}}A^{*}\mathbf{1}=\mathbf{1}^{\mathrm{T}}\bigl(|A|A^{-1}\bigr)\mathbf{1}}
      （$|A|\ne0$ 时），于是问题变成\textbf{先求 $A^{-1}$ 或直接算 $A^{*}$ 的一个二次型}；
\item 若系数不是全 $1$ 而是 $k_1,\dots,k_n$，则\textbf{逐位替换}即可，不必重新展开。
\end{enumerate}
【反哺点】9 讲把这一招按"用行列式 / 用矩阵 / 用特征值"三条路线讲，
但你只要记住一句话——\textbf{"代数余子式的系数 = 被换掉的那一行"}，
三条路线其实是同一件事的三种算法。考场上先写出 $\sum a_{ik}A_{ik}=|A|$，
再问"我要的系数是什么"，把系数填进被替换行，题就变成算一个行列式。
\end{fdeep}

\begin{fdeep}{$|A^{*}|=|A|^{\,n-1}$ 为什么成立，以及它的两个推论}
\textbf{推导只需一行}：在 $AA^{*}=|A|E$ 两边取行列式，用 $|kE|=k^{n}$，得
\fml{|A|\cdot|A^{*}|=\bigl||A|E\bigr|=|A|^{\,n}\;\Longrightarrow\;|A^{*}|=|A|^{\,n-1}\qquad(|A|\ne0)}
\textbf{$|A|=0$ 的情形}要单独说：此时 $r(A)\le n-1$。
若 $r(A)\le n-2$ 则 $A^{*}=O$，$|A^{*}|=0=|A|^{n-1}$ 成立；
若 $r(A)=n-1$ 则 $r(A^{*})=1$，$A^{*}$ 的各行成比例，仍得 $|A^{*}|=0=|A|^{n-1}$。
（这里正好复用了「$r(A^{*})$ 三分」那一条，两条结论互为支撑。）

\medskip
\textbf{两个常被忽略的推论}：
\begin{itemize}[leftmargin=2.2em,itemsep=0.22em]
\item $(A^{*})^{*}=|A|^{\,n-2}A$（$n\ge3$）。
      思路：对 $A^{*}$ 用同一公式，$(A^{*})^{*}=|A^{*}|(A^{*})^{-1}=|A|^{\,n-1}\cdot\dfrac{A}{|A|}=|A|^{\,n-2}A$；
      \textbf{注意 $n=2$ 时此式退化为 $(A^{*})^{*}=A$}，不要在 $n=2$ 的题里套错指数。
\item $|A^{*}|=|A|^{\,n-1}$ 的\textbf{反向用法}：见到 $|A^{*}|$ 与 $|A|$ 同时出现，
      可以直接把次数配平。例如由 $|A^{*}|=8$ 且 $n=3$ 立即得 $|A|^{2}=8$，$|A|=\pm2\sqrt2$ 需再定号。
\end{itemize}
【反哺点】考纲要求"会用伴随矩阵求逆矩阵"，而 $|A^{*}|$、$(A^{*})^{*}$ 这类量
正是填空题的常见落点。它们没有独立记忆的必要——全部挂在 $AA^{*}=|A|E$ 上现推即可，
但\textbf{$n=2$ 与 $n\ge3$ 的指数差别}必须记牢。
\end{fdeep}

\begin{fdeep}{所有元素都加上同一个数 $x$ 后的行列式与 $\sum\sum A_{ij}$}
\textbf{结论（例 2.55）}\quad 记 $A=(a_{ij})_{n\times n}$，$A_{ij}$ 是 $a_{ij}$ 的代数余子式，$J$ 为元素全为 $1$ 的 $n$ 阶方阵，则
\fml{\begin{vmatrix}
a_{11}+x & a_{12}+x & \cdots & a_{1n}+x\\
a_{21}+x & a_{22}+x & \cdots & a_{2n}+x\\
\vdots & \vdots & & \vdots\\
a_{n1}+x & a_{n2}+x & \cdots & a_{nn}+x
\end{vmatrix}=|A|+x\sum_{i=1}^{n}\sum_{j=1}^{n}A_{ij},}
即 $|A+xJ|=|A|+x\sum\limits_{i=1}^{n}\sum\limits_{j=1}^{n}A_{ij}$。特别地，$x=1$ 时
\fml{|A+J|=|A|+\sum_{i=1}^{n}\sum_{j=1}^{n}A_{ij}.}

\textbf{推导要点}\quad 只需三个动作，全程只用行列式的线性性质：
\fml{\begin{aligned}
|A+xJ|&=\begin{vmatrix}
a_{11}+x & a_{12}+x & \cdots & a_{1n}+x\\
a_{21} & a_{22} & \cdots & a_{2n}\\
\vdots & \vdots & & \vdots\\
a_{n1} & a_{n2} & \cdots & a_{nn}
\end{vmatrix}+\begin{vmatrix}
x & x & \cdots & x\\
a_{21} & a_{22} & \cdots & a_{2n}\\
\vdots & \vdots & & \vdots\\
a_{n1} & a_{n2} & \cdots & a_{nn}
\end{vmatrix}\\
&=\ \cdots\ =\ |A|+x\begin{vmatrix}
1 & 1 & \cdots & 1\\
a_{21} & a_{22} & \cdots & a_{2n}\\
\vdots & \vdots & & \vdots\\
a_{n1} & a_{n2} & \cdots & a_{nn}
\end{vmatrix}
\end{aligned}}
（对每一行都用"把第一行乘 $-1$ 加到其余各行，再把第一行拆项"，逐行做 $n$ 次），
最后那个行列式恰好是 $\sum\limits_{i=1}^{n}\sum\limits_{j=1}^{n}A_{ij}$。

\textbf{顺带得到的两件事}

（a）\textbf{余子式之和与 $x$ 无关}\quad 把 $D=|A+xJ|$ 中 $(i,j)$ 元的代数余子式记为 $D_{ij}$，则
\fml{\sum_{i=1}^{n}\sum_{j=1}^{n}D_{ij}=\sum_{i=1}^{n}\sum_{j=1}^{n}A_{ij}}
（因为 $|A|=D-x\sum\sum D_{ij}$，代回上式即得），故 $D$ 的所有元素的代数余子式之和与 $x$ 无关。

（b）\textbf{把 $x$ 取成 $-a_{11}$ 等具体值}，可把"含参行列式"化为"求 $\sum\sum A_{ij}$"；
反过来，$\sum\sum A_{ij}$ 又等于
\fml{\sum_{i=1}^{n}\sum_{j=1}^{n}A_{ij}=|A|\cdot\sum_{i=1}^{n}\sum_{j=1}^{n}\bigl(A^{-1}\bigr)_{ij}\qquad (|A|\neq0),}
即 $\sum\sum A_{ij}$ 与 $A^{-1}$ 的全部元素之和只差一个因子 $|A|$。

\textbf{【反哺点】}考纲内一旦出现"$\sum\sum A_{ij}$"或"所有元素同加一个数"，直接套 $|A+xJ|=|A|+x\sum\sum A_{ij}$：把已知行列式凑成 $A+xJ$ 的形状，未知的 $\sum\sum A_{ij}$ 就是一次多项式的一次项系数。
\end{fdeep}

\begin{fdeep}{行和、列和全为 $0$ 时所有代数余子式相等}
\textbf{结论（例 2.57）}\quad 设 $A=(a_{ij})_{n\times n}$ 满足
\fml{\begin{cases}
a_{i1}+a_{i2}+\cdots+a_{in}=0,\\[1mm]
a_{1j}+a_{2j}+\cdots+a_{nj}=0
\end{cases}\qquad(i,j=1,2,\cdots,n),}
（即 $A$ 的每一行、每一列元素之和都为 $0$）则 $A$ 的所有元素的代数余子式相等，即
\fml{A_{ij}=A_{11}\qquad(i,j=1,2,\cdots,n).}

\textbf{关键证明思路（$A_{ij}=A_{i1}$ 的一段）}\quad 对第一行的代数余子式，$j=2,3,\cdots,n$ 时
\fml{A_{1j}=(-1)^{1+j}\begin{vmatrix}
a_{21} & \cdots & a_{2,j-1} & a_{2,j+1} & \cdots & a_{2n}\\
a_{31} & \cdots & a_{3,j-1} & a_{3,j+1} & \cdots & a_{3n}\\
\vdots & & \vdots & \vdots & & \vdots\\
a_{n1} & \cdots & a_{n,j-1} & a_{n,j+1} & \cdots & a_{nn}
\end{vmatrix}.}
把第 $2$ 列至第 $n-1$ 列都加到第 $1$ 列，利用 $\sum\limits_{j=1}^{n}a_{ij}=0\ (i=2,3,\cdots,n)$ 得
\fml{A_{1j}=(-1)^{1+j}\begin{vmatrix}
-a_{2j} & a_{22} & \cdots & a_{2,j-1} & a_{2,j+1} & \cdots & a_{2n}\\
-a_{3j} & a_{32} & \cdots & a_{3,j-1} & a_{3,j+1} & \cdots & a_{3n}\\
\vdots & \vdots & & \vdots & \vdots & & \vdots\\
-a_{nj} & a_{n2} & \cdots & a_{n,j-1} & a_{n,j+1} & \cdots & a_{nn}
\end{vmatrix},}
再把第 $1$ 列的负号提出，并把第 $1$ 列依次与第 $2$ 列、第 $3$ 列、$\cdots$、第 $j-1$ 列互换（共 $j-2$ 次），得
\fml{A_{1j}=(-1)^{1+j}(-1)^{1}\cdot(-1)^{j-2}\begin{vmatrix}
a_{22} & \cdots & a_{2,j-1} & a_{2j} & a_{2,j+1} & \cdots & a_{2n}\\
a_{32} & \cdots & a_{3,j-1} & a_{3j} & a_{3,j+1} & \cdots & a_{3n}\\
\vdots & & \vdots & \vdots & \vdots & & \vdots\\
a_{n2} & \cdots & a_{n,j-1} & a_{nj} & a_{n,j+1} & \cdots & a_{nn}
\end{vmatrix}=A_{11}.}
对第 $i\ (i=2,3,\cdots,n)$ 行的代数余子式 $A_{ij}$ 重复上述过程，得 $A_{ij}=A_{i1}$；同理用 $a_{1j}+a_{2j}+\cdots+a_{nj}=0$ 可证 $A_{i1}=A_{11}$。故 $A_{ij}=A_{11}$。

\textbf{【反哺点】}遇到"行和与列和都为 $0$"的矩阵，$\sum\sum A_{ij}=n^{2}A_{11}$，只需算一个代数余子式；这也是判断"某矩阵的伴随矩阵是否为常矩阵（数量矩阵）"的最快入口。
\end{fdeep}

\begin{fdeep}{$A_{ij}=a_{ij}$ 的实矩阵：可逆性与 $|A|^{n-2}=1$}
\textbf{结论（例 2.60，上海师范大学 2002）}\quad 设 $a_{ij}\in\mathbb{R}$（实矩阵 $A=(a_{ij})_{n\times n}$），且 $A$ 中至少有一个元素 $a_{ij}\neq0$。若 $A$ 的一切元素 $a_{ij}$ 的代数余子式 $A_{ij}=a_{ij}$，则
\fml{D^{n-2}=1\qquad\bigl(D=|A|\bigr).}

\textbf{思路}\quad 取 $a_{ij}\neq0$，按其所在列展开：
\fml{D=\sum_{i=1}^{n}a_{ij}A_{ij}=\sum_{i=1}^{n}a_{ij}^{2}>0;}
又因 $A_{ij}=a_{ij}$ 即 $A^{*}=A^{T}$，由行列式乘法规则
\fml{D^{2}=DD^{T}=\begin{vmatrix}
a_{11} & a_{12} & \cdots & a_{1n}\\
a_{21} & a_{22} & \cdots & a_{2n}\\
\vdots & \vdots & & \vdots\\
a_{n1} & a_{n2} & \cdots & a_{nn}
\end{vmatrix}\begin{vmatrix}
A_{11} & A_{21} & \cdots & A_{n1}\\
A_{12} & A_{22} & \cdots & A_{n2}\\
\vdots & \vdots & & \vdots\\
A_{1n} & A_{2n} & \cdots & A_{nn}
\end{vmatrix}=\begin{vmatrix}
D & 0 & \cdots & 0\\
0 & D & \cdots & 0\\
\vdots & \vdots & & \vdots\\
0 & 0 & \cdots & D
\end{vmatrix}=D^{n},}
故 $D^{n-2}=1$。

\textbf{【反哺点】}考纲内凡是给"代数余子式与元素本身的关系"，都翻译成伴随矩阵语言 $A^{*}=A^{T}$（或 $A^{*}=f(A)$），再用 $AA^{*}=|A|E$ 一次算出 $|A|$；本题还顺带说明：实矩阵不必对称，只要 $A_{ij}=a_{ij}$ 就一定有 $|A|>0$。
\end{fdeep}

\fsec{深化三　与线性方程组的衔接}

\begin{fdeep}{用逆矩阵推导克拉默法则}
利用矩阵的逆，可以给出克拉默法则的另一种推导法.线性方程组
\fml{\begin{cases}
a_{11}x_1+a_{12}x_2+\cdots+a_{1n}x_n=b_1,\\
a_{21}x_1+a_{22}x_2+\cdots+a_{2n}x_n=b_2,\\
\cdots\cdots\cdots\cdots\\
a_{n1}x_1+a_{n2}x_2+\cdots+a_{nn}x_n=b_n
\end{cases}}
可以写成
\fml{AX=B.\qquad (6)}
如果 $|A|\ne 0$，那么 $A$ 可逆.用
\fml{X=A^{-1}B}
代入 (6)，得恒等式 $A(A^{-1}B)=B$，这就是说 $A^{-1}B$ 是一个解.

如果 $X=C$ 是 (6) 的一个解，那么由 $AC=B$ 得
\fml{A^{-1}(AC)=A^{-1}B,}
即
\fml{C=A^{-1}B.}
这就是说，解 $X=A^{-1}B$ 是唯一的.用 $A^{-1}$ 的公式 (4) 代入，乘出来就是克拉默法则中给出的公式.

【反哺点】服务考纲「线性方程组」"会用克拉默法则"：把"解方程组"统一成"求逆矩阵"，正是考纲「矩阵」与「线性方程组」两节的接口.这段还给出"存在性（代回验证 $A(A^{-1}B)=B$）$+$ 唯一性（两边左乘 $A^{-1}$）"的完整证明模板，凡是问"证明某矩阵表达式是唯一解"的题都可照此两步走.
\end{fdeep}